% ===========================================================================
\section{Limitations}
\label{sec:limitations}
% ===========================================================================

\textbf{The class split has no independent validation and should be read as the sector's
weakest output.} No nationally representative Pakistani vehicle-kilometre survey has been
identified for this study, and no on-road fuel-economy measurement either. The shares rest
on three inputs, none of them measured in the country: registration counts, assumed annual
kilometres, and the declared per-kilometre energy intensities that weight them. The largest
single result---two-wheelers at \SI{34.1}{\percent} of on-road \COtwo{} under the
reference allocation---rests on the least surveyed class in the fleet and on a judgement about how
much fuel a motorcycle burns.

This edition changed that result, which is itself a limitation worth stating: the previous
allocation weighted kilometres alone and put two-wheelers at \SI{42.3}{\percent}, above
heavy goods vehicles and buses combined. Correcting the weighting reversed that ordering.
A reader should conclude that the composition is not yet settled, not that it is now
settled differently. Anchoring confines the consequence to composition rather than total,
which is the reason the method was chosen, but it does not make the composition reliable.

\textbf{The off-road class is a residual as well as a category, and part of it belongs to a
different IPCC code.} The source workbook maps its catch-all "other vehicles" entry into
off-road wholesale, and the whole class is then modelled as \SI{100}{\percent} diesel. The
magnitude of the carve-out is therefore a genuine partition of the fuel total, while what
is inside it is soft. Whatever construction and industrial machinery the residual contains
belongs under 1A2\,gvii, not 1A4c; this inventory reports the whole class under the
mixture-declaring name \emph{Mixed off-road (1A4c/1A2)}, declares the misassignment in the
completeness table, and prices the reporting recode as a scenario, because splitting the
residual needs equipment stock data no source here provides. The cropland proxy that
allocates it provincially is right for tractors and wrong for whatever else the residual
contains.

\textbf{Retail-bought agricultural and captive diesel remain inside the road residual.}
Tube wells, threshers and captive generators fuelled at retail pumps draw on the Road
customer class, this inventory has not separated them, and they bias the road residual
high---conditional on the other deductions, since errors in the generator carve-out or
the off-road partition could run either way. The claim is framed rather than open-ended:
bulk agricultural HSD is observed zero in all ten OCAC editions; the ten-year road
share of national diesel (0.83--0.91) is a customer-class boundary sensitivity that
bounds the sales channel, not the contamination within it; and the retail non-road
scenarios price the within-channel claim at 2.5--10\% of the Road row
(Table~\ref{tab:boundary}). Resolving it needs a retail-channel end-use survey, which no
identified source provides. The ledger carries three claimants as \textsc{unresolved}
status rows---the third, "other retail stationary/non-road uses", because the named two
are not proven exhaustive.

\textbf{The provincial allocation now rides fuel sales, whose residual blur is transit.}
Fuel sold in a province is not fuel consumed there---freight crosses provincial
boundaries, so Sindh's \SI{19.45}{\percent} of road diesel partly reflects fuel bought at
Karachi and burned en route north. The shape is the reference year's own twelve area
pages---an earlier edition borrowed FY2024--25's here on the false premise that none were
published for the reference period, a borrowing since measured at no more than \num{0.45}
percentage points of error on any share. This is a far smaller distortion than the superseded
registration proxy's (which understated Sindh threefold; the comparison is published as a
diagnostic), but it is a real one. The release now ships the road geometry that a finer
allocation would start from---\texttt{outputs/pakemis\_road\_length\_adm2.csv} and a
\SI{0.05}{\degree} major-road density raster, extracted from a dated, hashed OSM
snapshot---as a diagnostic only: road length is not traffic, and moving fuel onto the
network would need traffic observations this inventory does not hold. CNG's provincial split still rides the fleet
model and the off-road class rides cropland (Table~\ref{tab:allocwalk}); only the
petroleum streams ride fuel sales.

\textbf{The CNG heating value is unpinned.} CNG is estimated from OGRA's metered sector
gas, but OGRA reports a volume and IPCC works in net terajoules: the conversion applies a
gross heating value and then the \num{0.90} gross-to-net factor for natural gas. The
gross-to-net step is a documented convention; the gross value is not---Pakistani pipeline gas
runs \SIrange{900}{1000}{} Btu/scf depending on field and RLNG share and nothing read here
pins it. Both ends are reported; the span is \SI{0.13}{\mega\tonne}, a tenth of the CNG
stream. This is the same unresolved quantity as the captive-power category's heating value,
and the two sectors are pinned to the same values by test so that one inventory cannot
assume the same gas two ways.

\textbf{The fleet model overstates CNG travel, and this sector has evidence rather than a
suspicion.} Section~\ref{sec:validation} divides the metered gas by the modelled CNG
vehicle-kilometres and gets \SI{0.91}{\mega\joule\per\kilo\metre}, below any plausible
value for a fleet seven-tenths three-wheelers by distance. The consequence is confined to
what rides vehicle-kilometres---the air diagnostic and the CNG streams' provincial
split---because the CNG emissions come from metered gas. It is nonetheless the clearest
evidence in this paper for the sector's central methodological choice.

\textbf{The air layer is an exploratory diagnostic, on factors without provenance.} The
per-kilometre factors are transcriptions of the Punjab emissions workbook, for which no
primary citation has been located---an earlier draft attributed them to the EMEP/EEA
guidebook, whose road-transport chapter has not in fact been acquired here. Evaporative
NMVOC and non-exhaust particulate are not estimated at all (both are named process
families in the gap register with the public chapters that would supply them), and
\SOtwo{} scales with an unreported fuel sulphur content. The layer is accordingly
reported in an appendix with no policy claims built on it.

\textbf{The time series is fuel-only, by design.} The ten-year fuel series
(Section~\ref{sec:series}) carries no generator back-cast, no off-road partition and no
class composition, because the inputs those refinements need exist only for the reference
year. The declared whole-Road-class basis keeps the series honest; it also means the
series is a magnitude record, not a structural one.

\textbf{Emission factors are Tier~1 defaults throughout.} They carry assumptions about
technology mix and combustion conditions derived largely from countries whose vehicle fleets
differ substantially from Pakistan's---most consequentially for the two-wheeler class, which
dominates here and is marginal in the fleets the defaults were derived from.

\textbf{The controls are member-supplied.} OCAC's own disclaimer states its data are
provided by members and not independently verified. The declared \SI{10}{\percent}
control uncertainty is anchored on the independent cross-check available---the railway
figures, \SI{3.8}{\percent} apart year-matched and \SI{11.2}{\percent} on the
conservative cross-year pairing---and the panel's internal closures, and it is a
judgement, not a fit.
